\section{二直線の交点では両 residue が二重零になる}

\begin{proposition}\label{prop:two-lines}
正規な実際の退化因子の被約特異支持に二本の直線が含まれ、
各直線の一般横断型が $\widetilde E_8$ であるとする。
両直線の交点における二つの modular residue の一次係数は共に零である。
交点に他の成分が入ることを許す。
\end{proposition}

\begin{proof}
両 slope を $a,b$ とする。一方だけが零である場合は
補題\ref{lem:zero-slope-nilpotent}に反するので、$a,b\ne0$ と仮定する。
前節の六つの spectrum を比較すると、$b/a=1$ または
$b/a=-1$、$h=h'=2$、$m=m'$ に限る。
最初の場合は tangent 二軸が同じ固有値 $a$ の独立固有ベクトルとなるが、
六つの spectrum における $a$ の重複度は一なので不可能。
従って $b=-a$, $h=h'=2$、$m\in\{0,1,2\}$ となる。
この比較は六つの明示的な四元多重集合の有理数比較であり、全場合を記録した。

両直線をアフィン $x,y$ 軸、残る座標を $z,w$ とする。
両 conormal bracket が交点で零なので、一次 Poisson tensor は
\[
 \pi_1=\ell(z,w)\partial_x\wedge\partial_y
\]
の形である。二つの直線の一定な Hessian covectors は、
$Z_1^*$ の単純固有値 $-a,a$ に対応する。
二軸を保つ線形座標変換により、それぞれを $dy,dx$ に取れる。
$m=0$ の場合、残る weight 零の二次元 block が冪零である可能性は除外しない。

\paragraph{$m=0$。}
$\ell=0$ なら $F$ の最低次数は四以上であり、交点で cubic scalar は零となる。
従って $m=0$ に反する。
$\ell\ne0$ なら、直線 conormal pencil が単純零を持つため
$\pi_2^{zw}$ の $xy$ 係数 $i$ は非零である。
三次 Pfaffian は $\ell\pi_2^{zw}$ であり、$i xy\ell$ は直線上で
一次の零しか持たない mixed Hessian を与える。他の一次行列成分はないので相殺されない。
これは Hessian の零点次数 $h=2$ に反する。

\paragraph{$m=1$。}
$a=3$ と規格化すると modular 線形化の固有値は $(3,-3,-1,1)$。
相異なる固有値に適合する線形座標を取る。
$[Z_1,\pi_1]=0$ は $\ell=0$ を与えるので、$\pi$ は二次から始まり、
$[Z_1,\pi_2]=0$ かつ $[\pi_2,\pi_2]=0$ である。
全ての重み零の二次成分は次の形となる。
\[
\begin{array}{ll}
P_{xy}=Axy+Bzw,&P_{xz}=Cxz+Dw^2,\\
P_{xw}=Exw,&P_{yz}=Gyz,\\
P_{yw}=Hyw+Iz^2,&P_{zw}=Jxy+Kzw.
\end{array}
\]
conormal pencil の単純零から $J\ne0$。
二つの cubic scalar の零点次数が一なので、Pfaffian の $xz^3,yw^3$ 係数が共に非零であり、
特に $D I\ne0$ である。
Jacobi の $xyz$ 成分における $xyz$ と $wz^2$ の係数は、それぞれ
\[
 BJ,\qquad B(C+G+K)-2DI.
\]
従って $B=0$ と $DI=0$ を得て矛盾する。

\paragraph{$m=2$。}
$a=3$ として固有値を $(3,-3,-2,2)$ に取る。
再び $\pi_1=0$、$[Z_1,\pi_2]=0$。
一定な Hessian と cubic の方向に適合した線形固有座標では、
二つの cubic scalar の二重零により、五次 Pfaffian の
$x^2z^3$ と $y^2w^3$ の係数が共に非零でなければならない。
これらはどちらも重み零である。

しかし $(3,-3,-2,2)$ に関して、重み零の二次ベクトル場は存在しない。
すなわち、二次単項式の重みは四つの座標重みのいずれとも一致しない。
従って $Z_2$ に重み零成分はない。
射影恒等式\eqref{eq:radial-cubic}と $E$ の重み零から、$\pi_3$ にも重み零成分はない。
$\pi_2$ は重み零、$\tr Z_1=0$ なので、混合 Pfaffian
$F_5$（$\pi_2\wedge\pi_3$ の係数）には重み零単項式がない。
先ほど必要とされた二つの非零係数と矛盾する。
ここでは線形座標変換だけを用いたので、射影恒等式の適用に問題はない。
\end{proof}

\paragraph{四直線の鎖と triple-tail。}
四直線の鎖では中央の直線上に二つの異なる modular 零点があり、
非零の $\PP^1$ 上の vector field の零因子次数二から各零点は単純である。
命題に反するので存在しない。
三重点と節点を持つ四直線でも、両点を結ぶ直線の節点側は単純零でなければならず、
同様に排除される。
conic--line--line の鎖も、中央の直線と末端直線の交点にこの議論を適用して排除できる。
